%% sections/conclusion.tex — Conclusion

We have presented a comprehensive numerical study of the Z2 quantum Mpemba effect
in the one-dimensional transverse-field Ising model and the anisotropic XY chain family.
The key contributions of this work are as follows.

Within the parameter range studied ($N \leq 18$, $\NA = 4$, $\theta \in [0.1, 1.9]$,
$\gf \in [0.1, 3.0]$), Z2 QME is a \emph{robust and pervasive} feature of
integrable free-fermion dynamics: QME crossings are observed for all
23 post-quench field values tested, spanning the ferromagnetic phase (no DQPTs
for $N=14$), the quantum critical point, and the paramagnetic phase (DQPTs present).
The near-complete occurrence rate — 69 field-pair combinations in the phase scan
and all detected crossings in the 671-pair statistical analysis — contrasts
with the expected absence of Z2 QME in random unitary circuits,
suggesting that integrability supports the effect through long-lived quasiparticle
coherence. Whether this extends to larger systems, other initial state families,
or non-integrable perturbations remains an open question.

We identify the \emph{initial entanglement asymmetry imbalance} $|\Delta\Delta\mathcal{S}|$
as the quantitative predictor of the crossing time $\tM$, with a Spearman correlation
$\rho = +0.90$ ($p \approx 3 \times 10^{-240}$): a larger initial EA gap produces a
later crossing, enabling $\tM$ to be tuned over nearly two orders of magnitude
(for pairs with $|\Delta\Delta\mathcal{S}| \geq 0.01$) purely through initial state preparation.
This initial-state origin of $\tM$ is made precise by the exact formula
$\DS_A(0;\theta) = \hbin((1 + \cos^{\NA}\!\theta)/2)$, which provides a
machine-precision-verified anchor and saturates the universal Z2 bound $\DS_A \leq \ln 2$.

We demonstrate, through rigorous statistical analysis of 671 pairs, that
\emph{dynamical quantum phase transitions do not govern QME crossing times}.
While the $\phi$ distribution is non-uniform (KS $p \approx 6 \times 10^{-25}$),
indicating that DQPT periodicity modulates the dynamics, the non-uniformity is
not centred on the DQPT location (binary test $p = 0.27$), and QME is present
even in the DQPT-free ferromagnetic phase.
DQPTs and QME thus belong to distinct categories of post-quench observables.

Finally, the crossing times converge smoothly to finite thermodynamic-limit values
(confirmed by $1/N$ extrapolation over $N = 8$--$18$) and are robust with respect
to subsystem size ($\NA = 2$--$6$) and model anisotropy ($\gamma = 0$--$1$),
establishing Z2 QME as a stable bulk phenomenon rather than a finite-size artefact.

These findings open a path toward a predictive theory of symmetry-restoration
dynamics in integrable systems, where the initial EA imbalance — directly encodable
through initial polarisation angles in cold-atom platforms — writes the timescale
of quantum Mpemba crossings.
